\section{Limitations, Monitoring, and Publication Decision}
\label{sec:limitations}

\subsection{Evidence limitations}

\begin{itemize}
  \item \textbf{Generated fixture.} No borrower or customer portfolio is evaluated. External
        validity, representativeness, and deployment drift are unknown.
  \item \textbf{Conditional inference.} Single-run confidence intervals and p-values condition on
        the generated rows and configured sample sizes. They exclude generator, scenario,
        seed-plan, model-selection, and deployment uncertainty.
  \item \textbf{Fixed seed suite.} Ten Gold seeds per scenario demonstrate completion under a
        declared suite; they are not a probabilistic sample of all future runs.
  \item \textbf{Mechanically imposed parity.} Intrinsic near-parity follows the equal-rate
        post-label policy. It is not causal evidence, an applicant-level counterfactual, or a
        finding that disparity was localized to decision labels.
  \item \textbf{Protected-group scope.} The fixture covers only declared gender and race encodings.
        It omits other protected bases, intersectional claims, encoding error, missingness
        mechanisms, and subgroup instability.
  \item \textbf{Race-policy split.} The headline evidence exporter uses the highest observed
        selection-rate group, while a separate configured-reference helper remains in the exact
        product source. The companion truthfully records white as configured and black as
        effective, but product-wide policy convergence is unresolved.
  \item \textbf{Utility weakness.} Synthetic-training ROC AUC ranges from
        \UtilitySyntheticAUCMin{} to \UtilitySyntheticAUCMax{} and is below the fixture baseline
        for every tested seed. Mean skill-normalised AUC retention is
        \UtilitySkillRetentionMean{} (\UtilitySkillRetentionMeanPct{}\%). This evidence does not
        establish train-on-synthetic predictive utility or support a production-utility claim.
  \item \textbf{Privacy scope.} Exact-row checks are narrow. Near-duplicate risk, membership
        inference, attribute inference, linkage, and formal privacy are not evaluated.
  \item \textbf{Decision-only metrics.} ECE is unevaluated and EO is non-informative. No score
        calibration, ranking quality, threshold optimization, adverse-action validity, or causal
        fairness analysis is established.
\end{itemize}

\subsection{Deployment monitoring pattern}

A customer-specific use would require a separately approved analysis plan. At minimum, control
owners should define the eligible population and decision point; validate protected-attribute
collection and lawful use; pre-specify group/reference policy, minimum counts, multiplicity,
missing-data handling and action thresholds; measure source and output drift; evaluate predictive
utility and calibration where scores exist; investigate exact and near duplicates; and retain
review, exception, remediation, and override records.

Monitoring should distinguish a statistical alert from a legal or business decision. An AIR point
below an internal line can trigger investigation without being labelled unlawful; an AIR above the
line cannot close an investigation where other evidence suggests harm. Confidence intervals,
effect sizes, data quality, operational context, and affected-person outcomes belong in the review.

\subsection{Remaining publication blockers}

The technical candidate can be reviewed, but the following items remain blockers to public
publication or stronger product claims:

\begin{enumerate}
  \item an authorized owner must approve the paper, exact PDF, companion ZIP, distribution route,
        and current legal wording;
  \item v5.0.1 is security-superseded because its lock pins \texttt{cryptography} 49.0.0 in the
        CVE-2026-69247 affected range. The exact v5.0.2 successor and fresh release evidence now
        exist, so this artifact's disposition must resolve to an explicitly archival v5.0.1
        characterization or retirement in favour of a separately rebuilt v5.0.2 paper; it cannot
        be relabelled or reused as successor evidence;
  \item the PDF and companion require a durable co-distribution location; this candidate has no
        public companion URL;
  \item a trusted external digest channel or verifiable release signature is required before
        describing the delivered companion as independently authenticated;
  \item the product's configured-reference and highest-rate race paths should be unified or
        explicitly versioned before claiming one product-wide race method;
  \item customer-evidence and promotion dispositions remain false; and
  \item production utility, broader protected-group coverage, near-duplicate privacy, calibration,
        and deployment security require additional evidence before related claims can be made.
\end{enumerate}

\subsection{Conclusion}

The exact v5.0.1 evidence supports a precise characterization: the amplification branch preserves
the fixture's gender disparity signal, the intrinsic parity-policy control operates as configured,
and the declared release Gold suite completes across 40 runs. The same evidence also exposes
important boundaries: conditional single-run inference, a split race-reference policy, incomplete
privacy and calibration surfaces, and materially variable downstream utility. Preserving both the
positive controls and the negative findings is the central governance value of this paper.
